\documentclass[../../main]{subfiles}

\ifSubfilesClassLoaded{%
    \zexternaldocument*{../1.Category-Theory/exported/aux/main}%
    \zexternaldocument*{../2.Homological-Algebra/exported/aux/main}%
    \zexternaldocument*{../3.Sheaves/exported/aux/main}%
    \zexternaldocument*{1.Cohomology-of-sheaves}%
    \zexternaldocument*{3.Examples-of-cohomology}%
}{}

\begin{document}
\chapter{de Rham Cohomology}

In this chapter, we bridge the gap between differential geometry and the abstract machinery of sheaf cohomology.
Historically, de Rham cohomology was developed to capture topological data about a smooth manifold using the analytical properties of differential forms.
Our primary goal here is to re-interpret this classical construction through a modern algebraic lens.

By treating spaces of differential forms as \hyperref[def:sheaf]{sheaves}, we will establish a fundamental isomorphism showing that classical de Rham cohomology is merely a realization of \hyperref[def:sheaf-cohomology]{sheaf cohomology} with coefficients in the \hyperref[def:constant-sheaf]{locally constant sheaf} \(\underline{\mathbb{R}}\).

Throughout this chapter we will work with \hyperref[def:sheaf]{sheaves of \(\mathbb{R}\)-modules}.
However, the same results follow \textit{mutatis mutandis} if we work only with \hyperref[def:sheaf]{sheaves of abelian groups}.
We will assume that all manifolds are real.
The complex case works analogously replacing \(\mathbb{R}\) with \(\mathbb{C}\) where appropiate.

\section{Differential forms and the de Rham complex}
We start this chapter by giving a quick recap of differential forms in smooth manifolds.
We refer the reader to \cite[Chapter 4]{guilleminDifferentialTopology1974a} and \cite[Chapter 14]{leeIntroductionSmoothManifolds2012} for more details.

To talk about differential forms globally and locally, we first need to assemble their point-wise linear-algebraic counterparts into a smooth geometric object.
This is accomplished by taking the exterior power of the cotangent bundle.
\begin{definition}\label{def:smooth-k-form-bundle}
  Let \(M\) be a smooth manifold of dimension \(m\).

  The \textbf{bundle of \(k\)-forms} \(\Lambda^{k}(\mathrm{T}^{\,\ast}M)\) is the \hyperref[def:smooth-vector-bundle]{smooth vector bundle} defined with:
  \begin{itemize}
    \item \textit{Total space}: \(\Lambda^{k}(\mathrm{T}^{\,\ast}M) = \coprod_{p \in M} \Lambda^{k}(\mathrm{T}_{p}^{\,\ast}M)\).
    \item \textit{Projection map}: \(\pi: \Lambda^{k}(\mathrm{T}^{\,\ast}M) \to M,\; \pi(p, \omega) = p\).
    \item \textit{Fiber space}: For every \(p \in M\) the fiber \(E_{p} = \pi^{-1}(\{p\}) = \Lambda^{k}(\mathrm{T}_{p}^{\,\ast}M)\) is a real vector space of dimension \(\binom{m}{k}\).
  \end{itemize}
  To define the local trivialization, consider a local chart \((U, \{x^{i}\}_{i=1}^{m})\) of \(M\).
  Then the differentials \(\{\mathrm{d}x^{i_{1}} \wedge \cdots \wedge \mathrm{d}x^{i_{k}}\}_{1 \leq i_{1} < \cdots < i_{k} \leq m}\) form a basis for each fiber with base point in \(U\) and the map \(\phi^{-1}: \pi^{-1}(U) \to U \times \mathbb{R}^{{m \choose k}}\) is given by sending a form \(\omega\) to its coefficients in this base.
\end{definition}

\begin{notation}
  To simplify working with these local bases, we adopt the standard multi-index notation for the remainder of the text.
  Let \(I = (i_{1}, \ldots, i_{k})\) be a multi-index.
  We will write
  \begin{equation*}
    \mathrm{d}x^{I} = \mathrm{d}x^{i_{1}} \wedge \cdots \mathrm{d}x^{i_{k}}
  \end{equation*}
  In this way the basis of the local trivialization is \(\{\mathrm{d}x^{I}\}_{I}\).
\end{notation}

With the underlying vector bundle established, we can now define differential forms dynamically as smoothly varying choices of tensors across the manifold.
In the language of sheaf theory, this naturally corresponds to taking the sheaf of sections of our bundle.

\begin{definition}[Smooth \(k\)-form] \label{def:smooth-k-form}
  Let \(M\) be a smooth manifold of dimension \(m\).

  The \textbf{\hyperref[def:sheaf]{sheaf} of \(k\)-forms} \(\Omega^{k}\) over \(M\) is the \hyperref[def:sheaf-of-sections]{sheaf of sections} of the \hyperref[def:smooth-k-form-bundle]{bundle of \(k\)-forms} \(\Lambda^{k}(\mathrm{T}^{\,\ast}M)\).

  More concretely, a \(k\)-form over an open set \(U\) is a map \(\omega: U \to \Lambda^{k}(\mathrm{T}^{\,\ast}M)\) such that \(\pi \circ \omega = \id_{U}\).
\end{definition}

\begin{example}
  Let \(M\) be a smooth manifold of dimension \(m\).

  Note that \(\Lambda^{0}(\mathrm{T}^{\,\ast}M) = M \times \mathbb{R}\) is the \hyperref[ex:trivial-topological-vector-bundle]{trivial vector bundle}.
  Therefore a \hyperref[def:smooth-k-form]{\(0\)-form} is a map \(\omega: M \to M \times \mathbb{R}\) such that \(\pi \circ \omega = \id_{M}\), i.e: \(\omega\) is a smooth function.

  This means that \(\Omega^{0}(M) \cong \mathcal{C}^{\infty}(M)\).
\end{example}

The core driving force behind de Rham theory is the existence of a differential operator that links these sheaves together by increasing the degree of the forms.
Locally, this operator behaves like a directional extension of the standard total derivative.

\begin{definition}[Exterior derivative] \label{def:exterior-derivative-k-form}
  Let \(\omega\) be a \hyperref[def:smooth-k-form]{\(k\)-form} over a smooth manifold \(M\).
  If \((U, \{x^{i}\}_{i=1}^{m})\) is a local chart then there exist functions \(\alpha_{I} \in \mathcal{C}^{\infty}(M)\) such that
  \begin{equation*}
    \omega = \sum_{I} \alpha_{I} \mathrm{d}x^{I}
  \end{equation*}
  where \(I = (i_{1}, \ldots, i_{k})\) is a multi-index and \(\mathrm{d}x^{I} = \mathrm{d}x^{i_{1}} \wedge \cdots \mathrm{d}x^{i_{k}}\).

  We can define the \textbf{exterior derivative} of \(\omega\) as the \(k+1\)-form
  \begin{equation*}
    \mathrm{d}\omega = \sum_{I} \mathrm{d}\alpha_{I} \wedge \mathrm{d}x^{I} = \sum_{I} \left( \sum_{j=1}^{m} \frac{\partial \alpha_{I}}{\partial x^{j}} \mathrm{d}x^{j} \right) \wedge \mathrm{d}x^{I}
  \end{equation*}
\end{definition}

For these operators to chain together into a cochain complex, the composition of two successive derivatives must vanish.
This works thanks to the symmetry of partial derivatives and the antisymmetry of the wedge product.
\begin{lemma} \label{lemma:double-exterior-derivative-is-zero}
  Let \(M\) be a smooth manifold of dimension \(m\).

  Then for any \hyperref[def:smooth-k-form]{\(k\)-form} the \hyperref[def:exterior-derivative-k-form]{exterior derivative} verifies \(\mathrm{d}^{2} \omega = 0\), i.e: the sequence
  \begin{center}
    \begin{tikzcd}
      \Omega^{0}(M) \arrow[r, "\mathrm{d}"] & \Omega^{1}(M) \arrow[r, "\mathrm{d}"] & \Omega^{2}(M) \arrow[r, "\mathrm{d}"] & \cdots
    \end{tikzcd}
  \end{center}
  is a \hyperref[def:complex]{complex} of \hyperref[def:sheaf-of-modules]{sheaves of \(\mathbb{R}\)-modules}.
\end{lemma}
\begin{proof}
  It suffices to show that for any smooth function \(f \in \mathcal{C}^{\infty}(M) \cong \Omega^{0}(M)\) we have \(\mathrm{d}^{2}f = 0\).
  Indeed, if \(\mathrm{d}^{2}f = 0\) then for any basic form \(f\mathrm{d}x^{I}\) we have
  \begin{equation*}
    \mathrm{d}(f\mathrm{d}x^{I}) = \mathrm{d}(\mathrm{d}f \wedge \mathrm{d}x^{I}) = \mathrm{d}^{2}f \wedge \mathrm{d}x^{I} - \mathrm{d}f \wedge \mathrm{d}^{2}x^{I} = 0
  \end{equation*}

  Let \(f \in \mathcal{C}^{\infty}(M)\) be a smooth function.
  Its \hyperref[def:exterior-derivative-k-form]{exterior derivative} is
  \begin{equation*}
    \mathrm{d}f = \sum_{i = 1}^{m} \frac{\partial f}{\partial x^{i}} \mathrm{d}x^{i}
  \end{equation*}
  Differentiating again we get
  \begin{equation*}
    \mathrm{d}^{2}f = \sum_{j=1}^{m} \frac{\partial}{\partial x^{j}} \left( \sum_{i=1}^{m}  \frac{\partial f}{\partial x^{i}} \right) \wedge \mathrm{d}x^{j} = \sum_{i,j=1}^{m} \frac{\partial^{2} f}{\partial x^{i} \partial x^{j}} \mathrm{d}x^{i} \wedge \mathrm{d}x^{j}
  \end{equation*}
  By Clairaut's theorem we have that \( \frac{\partial^{2} f}{\partial x^{i} \partial x^{j}} =  \frac{\partial^{2} f}{\partial x^{j} \partial x^{i}}\) and by the antisymmetry of the wedge product \(\mathrm{d}x^{i} \wedge \mathrm{d}x^{j} = - \mathrm{d}x^{j} \wedge \mathrm{d}x^{j}\).
  Therefore all the terms cancel leaving
  \begin{equation*}
    \mathrm{d}^{2} f = 0 \qedhere
  \end{equation*}
\end{proof}

\begin{example}\label{ex:finiteness-de-Rham}
  Let \(M\) be a smooth manifold of dimension \(m\).

  Suppose that \(k > m\).
  Since the fibers of the \hyperref[def:smooth-k-form-bundle]{bundle of \(k\)-forms} \(\Lambda^{k}(\mathrm{T}^{\,\ast}M)\) have dimension \(\binom{m}{k}\) this means that at least one index \(i\) in the multi-index \(I\) will be repeated.
  Therefore since \(\mathrm{d}x^{i} \wedge \mathrm{d}x^{i} = 0\) we have that \(\mathrm{d}x^{I} = 0\), i.e: \(\Lambda^{k}(\mathrm{T}^{\,\ast}M) = 0\).

  In particular \(\Omega^{k}(M) = 0\) for all \(k > m\).
\end{example}

We can now introduce the de Rham complex and its associated cohomology.

\begin{definition}[de Rham complex] \label{def:de-rham-complex}
  Let \(M\) be a smooth manifold of dimension \(m\).
  The \textbf{de Rham complex} of \(M\) is the \hyperref[def:complex]{complex} of \hyperref[def:sheaf-of-modules]{sheaves of \(\mathbb{R}\)-modules}
  \begin{center}
    \begin{tikzcd}
      \Omega^{\bullet}(M) \coloneq \Omega^{0}(M) \arrow[r, "\mathrm{d}"] & \Omega^{1}(M) \arrow[r, "\mathrm{d}"] & \cdots \arrow[r, "\mathrm{d}"] & \Omega^{m}(M) \arrow[r, "\mathrm{d}"] & 0
    \end{tikzcd}
  \end{center}
  where \(\Omega^{k}(M)\) is the \hyperref[def:smooth-k-form]{sheaf of \(k\)-forms} over \(M\) and \(\mathrm{d}\) is the \hyperref[def:exterior-derivative-k-form]{exterior derivative}.
\end{definition}

\begin{definition}
  Let \(M\) be a smooth manifold.
  We define the \textbf{de Rham cohomology} of \(M\) as the \hyperref[def:cohomology]{cohomology} of the \hyperref[def:de-rham-complex]{de Rham complex} and denote it by \(\operatorname{H}_{\operatorname{dR}}^{\bullet}(M, \mathbb{R})\).
\end{definition}


\section{Fine sheaves}

It is known that all smooth manifolds (with or without boundary) admit a partition of unity (see \cite[Theorem 2.23]{leeIntroductionSmoothManifolds2012}).
In this section, we are going to see how we can use this analytical tool to compute \hyperref[def:sheaf-cohomology]{sheaf cohomology}.

First we need to extend the concept of a partition of unity to a \hyperref[def:sheaf]{sheaf}.
\begin{definition}[Partition of unity]\label{def:partition-of-unity}
  Let \(\sheaf{F}\) be a \hyperref[def:sheaf]{sheaf of \(\mathbf{k}\)-modules} over \(X\).
  A \textbf{partition of unity} of \(\sheaf{F}\) subordinate to a locally finite open covering \(\mathcal{U} = \{U_{i}\}_{i \in \mathcal{I}}\) of \(X\) is a collection of endomorphisms 
  \begin{equation*}
    \{ f_{i}: \sheaf{F} \to \sheaf{F} \}_{i \in \mathcal{I}}
  \end{equation*}
  satisfying
  \begin{enumerate}
    \item For all \(i \in \mathcal{I}\)
          \begin{equation*}
            \operatorname{supp} f_{i} \subset U_{i}
          \end{equation*}
    \item
          \begin{equation*}
            \sum_{i \in \mathcal{I}} f_{i} = \id_{\sheaf{F}}
          \end{equation*}
        Note that this sum is well defined since for any point \(x \in X\) there exists an open neighbourhood where all but finitely many of the maps \(f_{i}\) are the zero maps, i.e: the sum is locally finite.
  \end{enumerate}
\end{definition}

Now we can introduce the objects of study of this section: fine sheaves.
Roughly speaking these a \hyperref[def:sheaf]{sheaves} such that whose structure allows them to be split and glued together using these localized operators.
\begin{definition}[Fine sheaf]\label{def:fine-sheaf}
  A \hyperref[def:sheaf]{sheaf of \(\mathbf{k}\)-modules} \(\sheaf{F}\) is \textbf{fine} if for every open cover \(\mathscr{U} = \{U_{i}\}_{i \in \mathcal{I}}\) of \(X\) there exists a \hyperref[def:partition-of-unity]{partition of unity} \(\{f_{i}\}_{i \in \mathcal{I}}\) subordinate to this covering.
\end{definition}

\begin{proposition} \label{prop:fine-are-acyclic}
  \hyperref[def:fine-sheaf]{Fine sheaves} are \hyperref[def:F-acyclic-object]{\(\Gamma(X, -)\) acyclic}.
\end{proposition}
\begin{proof}
  Let \(\sheaf{F}\) be a \hyperref[def:fine-sheaf]{fine sheaf}.
  Using the \hyperref[def:godement-resolution]{Godement resolution} we have that \(\sheaf{F}\) admits a \hyperref[def:flasque-sheaf]{flasque} resolution \(\sheaf{F} \to \sheaf{I}^{\bullet}\).

  Using \autoref{prop:flasque-are-acyclic} by \autoref{thm:right-derived-functor-using-acyclics} we have that
  \begin{equation*}
    \cohomology{k}(X; \sheaf{F}) = \quot{\ker\left( \Gamma(X, \sheaf{I}^{k}) \to \Gamma(X, \sheaf{I}^{k+1}) \right)}{\operatorname{Im} \left( \Gamma(X, \sheaf{I}^{k-1}) \to \Gamma(X, \sheaf{I}^{k}) \right)}
  \end{equation*}
  Let \(\alpha \in \ker\left( \Gamma(X, \sheaf{I}^{k}) \to \Gamma(X, \sheaf{I}^{k+1}) \right)\) with \(k > 0\).
  Since \(\sheaf{I}^{\bullet}\) is locally exact in degree greater than 0 we have that \(\alpha\) comes from \(\sheaf{I}^{k-1}\), i.e: there exists an open cover \(\mathscr{U} = \{U_{i}\}_{i \in \mathcal{I}}\) of \(X\) such that \(\alpha|_{U_{i}} = \d\beta_{i}\) with \(\beta_{i} \in \sheaf{I}^{k-1}(U_{i})\).

  Consider a \hyperref[def:partition-of-unity]{partition of unity} \(\{f_{i}\}_{i \in \mathcal{I}}\) subordinate to \(\mathscr{U}\) and let
  \begin{equation*}
    \beta \coloneq \sum_{i \in \mathcal{I}} f_{i}\beta_{i}
  \end{equation*}
  where the sum is locally finite.
  Note that \(f_{i}\beta_{i}\) is a section of \(\sheaf{I}^{k-1}(X)\) with values in \(\sheaf{I}_{p}^{k-1}(X)\):
  \begin{equation*}
    (f_{i}\beta_{i})_{p} = \begin{cases}
      0 & p \notin U_{i} \\
      f_{i}\beta_{i} & p \in U_{i}
    \end{cases}
  \end{equation*}
  Since \(f_{i}\) has support in \(U_{i}\) this defines a section in \(\sheaf{I}^{k-1}(U_{i})\).

  We have that \(\d\beta = \alpha\) since \(\alpha = \sum_{i \in \mathcal{I}} f_{i}\alpha|_{U_{i}}\).
  Therefore \(\ker\left( \Gamma(X, \sheaf{I}^{k}) \to \Gamma(X, \sheaf{I}^{k+1}) \right) = \im \left( \Gamma(X, \sheaf{I}^{k-1}) \to \Gamma(X, \sheaf{I}^{k}) \right)\), which means
  \begin{equation*}
    \forall k > 0 \quad \cohomology{k}(X, \sheaf{F}) = 0 \qedhere
  \end{equation*}
\end{proof}

An example of \hyperref[def:fine-sheaf]{fine sheaves} is one we have just studied: \hyperref[def:smooth-k-form]{sheaves of \(k\)-forms}.
\begin{lemma} \label{lemma:smooth-k-form-fine}
  Let \(M\) be a smooth manifold.
  Then for all \(k\) the \hyperref[def:smooth-k-form]{sheaf of \(k\)-forms} \(\Omega^{k}\) is \hyperref[def:fine-sheaf]{fine}.
\end{lemma}
\begin{proof}
  Let \(\mathscr{U} = \{U_{i}\}_{i \in \mathcal{I}}\) be an arbitrary open cover of \(M\).

  Since \(M\) is paracompact there exists a partition of unity \(\{f_{i}\}_{i \in \mathcal{I}} \subset \mathcal{C}^{\infty}(M)\) subordinate to \(\mathscr{U}\).
  This means that
  \begin{align*}
    \forall i \in \mathcal{I} \quad& \operatorname{supp} f_{i} \subseteq U_{i} \\
    \forall p \in M \quad& \sum_{i \in \mathcal{I}} f_{i}(p) = 1
  \end{align*}
  and this last sum is finite for each \(p\).

  For all indices \(i \in \mathcal{I}\) consider the map \(\eta_{i}\colon \Omega^{k} \to \Omega^{k}\) defined as, for all \(U \subseteq M\) open
  \begin{align*}
    \eta_{i}(U):\,& \Omega^{k}(U) \to \Omega^{k}(U) \\
            & \omega \mapsto f_{i} \cdot \omega
  \end{align*}
  Note that since \(\operatorname{supp} f_{i} \subseteq U_{i}\) we have \(\operatorname{supp} \eta_{i} \subseteq U_{i}\).
  Furthermore for any \hyperref[def:smooth-k-form]{\(k\)-form} \(\omega\)
  \begin{equation*}
    \sum_{i \in \mathcal{I}} \eta_{i}(\omega) = \sum_{i \in \mathcal{I}} f_{i} \cdot \omega = \left( \sum_{i \in \mathcal{I}} f_{i} \right) \cdot \omega = 1 \cdot \omega = \omega
  \end{equation*}
  Therefore \(\sum_{i \in \mathcal{I}} \eta_{i} = \id_{\Omega^{k}}\), making \(\{\eta_{i}\}_{i \in \mathcal{I}}\) a partition of unity subordinate to \(\mathscr{U}\).
\end{proof}

Therefore if we can prove that there exists a \hyperref[def:sheaf]{sheaf of \(\mathbb{R}\)-modules} \(\sheaf{F} \in \Mod(\underline{\mathbb{R}}_{M})\) such that \(\Omega^{\bullet}\) is a \hyperref[def:resolution]{resolution} of \(\sheaf{F}\) we would have that \hyperref[def:sheaf-cohomology]{sheaf cohomology with values in \(\sheaf{F}\)} and \hyperref[def:de-rham-cohomology]{de Rham cohomology} coincide.

As we are going to see, this is almost a direct consequence of the classical Poincaré lemma of multi-variable calculus if we consider \(\sheaf{F} = \underline{\mathbb{R}}_{M}\).
\begin{lemma}[{Poincaré lemma --- \cite[Theorem 11.49]{leeIntroductionSmoothManifolds2012}}]\label{lemma:poincare}
  Let \(U \subseteq \mathbb{R}^{n}\) be a contractible open set and \(k \geq 1\).
  If \(\omega \in \Omega^{k}(U)\) is a \hyperref[def:smooth-k-form]{\(k\)-form} such that \(\mathrm{d}\omega = 0\) then there exists a \hyperref[def:smooth-k-form]{\((k-1)\)-form} \(\eta \in \Omega^{k-1}(U)\) such that \(\mathrm{d}\eta = \omega\).
\end{lemma}

First, let's check that \(\Omega^{\bullet}\) is an \hyperref[def:exact-sequence]{exact complex} of \hyperref[def:sheaf]{sheaves}.
\begin{lemma} \label{lemma:de-rham-complex-exact}
  The \hyperref[def:complex]{complex} of \hyperref[def:sheaf]{sheaves} \(0 \to \Omega^{0} \to \Omega^{1} \to \cdots\) is \hyperref[def:exact-sequence]{exact}.
\end{lemma}
\begin{proof}
  For \(k \geq 1\) consider the relevant part of the complex
  \begin{center}
    \begin{tikzcd}
      \Omega^{k-1} \arrow[r, "\mathrm{d}^{k-1}"] & \Omega^{k} \arrow[r, "\mathrm{d}^{k}"] & \Omega^{k+1}
    \end{tikzcd}
  \end{center}
  We need to show that for any \(p \in M\) the complex
  \begin{center}
    \begin{tikzcd}
      \left(\Omega^{k-1}\right)_{p} \arrow[r, "\mathrm{d}_{p}^{k-1}"] & \left(\Omega^{k}\right)_{p} \arrow[r, "\mathrm{d}_{p}^{k}"] & \left(\Omega^{k+1}\right)_{p}
    \end{tikzcd}
  \end{center}
  is \hyperref[def:exact-complex]{exact}, i.e: \(\im \mathrm{d}_{p}^{k-1} = \ker \mathrm{d}_{p}^{k}\).

  For the \([\subseteq]\) inclusion, let \(\mathrm{d}\eta \in \im \mathrm{d}^{k-1} \subseteq \Omega^{k}\). By \autoref{lemma:double-exterior-derivative-is-zero} we have that \(\mathrm{d}(\mathrm{d}\eta) = 0\), i.e: \(\mathrm{d}\eta \in \ker \mathrm{d}^{k}\).

  For the \([\supseteq]\) inclusion, let \([\omega]_{p} \in \ker \mathrm{d}_{p}^{k}\), i.e: \([\omega]_{p} \in \left( \Omega^{k} \right)_{p}\) verifies \(\mathrm{d}_{p}^{k}[\omega]_{p} = 0\).
  By definition there exists an open neighbourhood \(U\) of \(p\) and a \hyperref[def:smooth-k-form]{\(k\)-form} \(\omega\) such that \(\mathrm{d}^{k}\omega|_{U} = 0\).

  Since \(M\) is a smooth manifold of dimension \(m\) we can choose a local chart \((V, \psi)\) around \(p\) such that \(V \subseteq U\) and \(\psi(V)\) is an open ball in \(\mathbb{R}^{m}\).
  Since open balls are contractible, by \hyperref[lemma:poincare]{Poincaré's lemma} there exists a \hyperref[def:smooth-k-form]{\((k-1)\)-form} \(\eta \in \Omega^{k-1}(V)\) such that \(\mathrm{d}^{k-1}\eta = \omega|_{V}\).

  Therefore there exists \([\eta]_{p} \in \left( \Omega^{k-1} \right)_{p}\) such that \(\mathrm{d}_{p}^{k-1}[\eta]_{p} = [\omega]_{p}\), i.e: \([\omega]_{p} \in \im d_{p}^{k-1}\).
\end{proof}

Having seen this, we can finally state and prove the main theorem of the chapter.
\begin{theorem}
  Let \(M\) be a smooth manifold.
  Then
  \begin{equation*}
     \cohomology{\bullet}(M, \underline{\mathbb{R}}_{M}) \cong \cohomology{\bullet}_{\operatorname{dR}}(M, \mathbb{R})
  \end{equation*}
\end{theorem}
\begin{proof}
  Consider the resolution
  \begin{center}
    \begin{tikzcd}
      0 \arrow[r] & \underline{\mathbb{R}} \arrow[r] & \Omega^{\bullet}
    \end{tikzcd}
  \end{center}
  where \(\Omega^{\bullet}\) is the \hyperref[def:de-rham-complex]{de Rham complex} of \(M\).

  By \autoref{lemma:de-rham-complex-exact} we know that \(\Omega^{\bullet}\) is \hyperref[def:exact-complex]{exact}, therefore we only need to show that this resolution is \hyperref[def:exact-complex]{exact} at \(\Omega^{0}\), i.e: that for all \(p \in M\)
  \begin{center}
    \begin{tikzcd}
      0 \arrow[r] & \underline{\mathbb{R}}_{p} \arrow[r, "\iota_{p}"] & \Omega_{p}^{0} \arrow[r, "\mathrm{d}_{p}"] & \Omega_{p}^{1} \arrow[r] & 0
    \end{tikzcd}
  \end{center}
  is a \hyperref[def:short-exact-sequence]{short exact sequence}, i.e: \(\im \iota_{p} = \ker \mathrm{d}_{p}\).

  \begin{itemize}
    \item[\(\subseteq\)] Let \([c]_{p} \in \underline{\mathbb{R}}_{p}\). By definition this is a locally constant function \(c: U \to \mathbb{R}\) on some open neighbourhood \(U\) of \(p\).

          Since \(c\) is locally constant we have that \(\mathrm{d}c|_{U} = 0\), i.e \(\mathrm{d}_{p}[c]_{p} = 0\).
    \item[\(\supseteq\)] Let \([f]_{p} \in \Omega_{p}^{0}\) such that \(\mathrm{d}_{p}[f]_{p} = 0\).
          This means that there exists an open neighbourhood \(U\) of \(p\) and a function \(f \in \Omega^{0}(V) \cong \mathcal{C}^{\infty}(M)\) such that \(\mathrm{d}f|_{U} = 0\).

          Therefore \(f\) is locally constant in each of the connected components of \(U\), i.e: \([f]_{p} \in \underline{\mathbb{R}}_{p}\).
  \end{itemize}

  Therefore since by \autoref{lemma:smooth-k-form-fine} the resolution is made of \hyperref[def:fine-sheaf]{fine sheaves} and \hyperref[prop:fine-are-acyclic]{fine sheaves are \(\Gamma(X, -)\) acyclic} by definition of \hyperref[def:sheaf-cohomology]{sheaf cohomology}
  \begin{equation*}
      \cohomology{\bullet}(M, \underline{\mathbb{R}}_{M}) \cong \cohomology{\bullet}_{\operatorname{dR}}(M, \mathbb{R}) \qedhere
  \end{equation*}
\end{proof}

As a consequence of \autoref{ex:finiteness-de-Rham} we get the following finiteness result for \hyperref[def:sheaf-cohomology]{sheaf cohomology}.
\begin{corollary}
  Let \(M\) be a smooth manifold of dimension \(m\).
  Then for all \(k > m\)
  \begin{equation*}
    \cohomology{k}(M, \underline{\mathbb{R}}_{M}) = 0
  \end{equation*}
\end{corollary}

To end this chapter we revisit the computation of the cohomology of the circle \(\mathbb{S}^{1}\) to verify that our results indeed hold.

\begin{example}[The cohomology of \(\mathbb{S}^{1}\)]
  Consider the circle \(\mathbb{S}^{1}\) as a smooth manifold of dimension 1.

  We are going to compute its \hyperref[def:sheaf-cohomology]{sheaf cohomology} with coefficients in \(\underline{\mathbb{R}}\), \(\cohomology{\bullet}(\mathbb{S}^{1}, \underline{\mathbb{R}})\), using the \hyperref[def:de-rham-complex]{de Rham complex}
  \begin{center}
    \begin{tikzcd}
      \Omega^{0}(\mathbb{S}^{1}) \arrow[r, "\d^{0}"] & \Omega^{1}(\mathbb{S}^{1}) \arrow[r, "\d^{1}"] & 0
    \end{tikzcd}
  \end{center}

  On the one hand
  \begin{equation*}
    \cohomology{0}(\mathbb{S}^{1}, \underline{\mathbb{R}}_{M}) = \ker \d^{0} = \{f \in \Omega^{0}(\mathbb{S}^{1}) \cong \mathcal{C}^{\infty}(\mathbb{S}^{1}) \mid \d f = 0\}
  \end{equation*}
  Since \(\mathbb{S}^{1}\) is connected, the only functions with null derivative are the constant functions.
  Therefore
  \begin{equation*}
    \cohomology{0}(\mathbb{S}^{1}, \underline{\mathbb{R}}_{M}) = \mathbb{R}
  \end{equation*}

  On the other
  \begin{equation*}
    \cohomology{1}(\mathbb{S}^{1}, \underline{\mathbb{R}}_{M}) = \quot{\ker \d^{1}}{\im \d^{0}} = \quot{\{\omega \in \Omega^{1}(\mathbb{S}^{1}) \mid \d \omega = 0\}}{\{\d f \mid f \in \Omega^{0}(\mathbb{S}^{1}) \cong \mathcal{C}^{\infty}(\mathbb{S}^{1})\}}
  \end{equation*}

  Note that since \(\Omega^{2}(M) = 0\) all \hyperref[def:smooth-k-form]{\(1\)-forms} \(\omega \in \Omega^{1}(\mathbb{S}^{1})\) verify \(\mathrm{d}\omega = 0\).
  Therefore
  \begin{equation*}
    \{\omega \in \Omega^{1}(\mathbb{S}^{1}) \mid \mathrm{d}\omega = 0\} = \Omega^{1}(\mathbb{S}^{1})
  \end{equation*}
  Consider now the \textit{integration} map
  \begin{align*}
    I:\,& \Omega^{1}(\mathbb{S}^{1}) \to \mathbb{R} \\
        & \omega \mapsto \int_{\mathbb{S}^{1}} \omega 
  \end{align*}
  Let \(\mathrm{d}\theta\) be the differential of the angle.
  Since \(I(\mathrm{d}\theta) = 2\pi\) we have that \(I\) is surjective.
  Furthermore by Stoke's theorem we have that \(I(\omega) = 0\) if and only if \(w = \mathrm{d}f\) for some \(f \in \mathcal{C}^{\infty}(\mathbb{S}^{1})\).

  Therefore by the first isomorphism theorem we have that
  \begin{equation*}
    \cohomology{1}(\mathbb{S}^{1}, \underline{\mathbb{R}}_{M}) = \quot{\Omega^{1}(\mathbb{S}^{1})}{\ker I} \cong \im I = \mathbb{R}
  \end{equation*}

  In summary, we have that
  \begin{align*}
    \cohomology{k}(\mathbb{S}^{1}, \underline{\mathbb{R}}_{M}) = \begin{cases}
      \mathbb{R} & k \in \{0, 1\} \\
      0 & k > 1
    \end{cases}
  \end{align*}
  which agrees with the \hyperref[ex:cohomology-S1]{our previous computation using only sheaf cohomology}.
\end{example}

\end{document}
